% BHU PYQ -- Professional Exam Template
\documentclass[11pt,a4paper]{article}

\usepackage[a4paper,top=2.5cm,bottom=2.5cm,left=2.8cm,right=2.8cm,headheight=14pt]{geometry}
\usepackage{amsmath,amssymb}
\usepackage[T1]{fontenc}
\usepackage[utf8]{inputenc}
\usepackage{lmodern}
\usepackage{microtype}
\usepackage{fancyhdr}
\usepackage{enumitem}
\usepackage{hyperref}
\usepackage{lastpage}
\usepackage{parskip}

\hypersetup{colorlinks=true,urlcolor=black,linkcolor=black,
  pdftitle={B.Sc. (Hons.) Botany Sem-VI BOB-602 -- BHU PYQ}}

\pagestyle{fancy}
\fancyhf{}
\renewcommand{\headrulewidth}{0.4pt}
\renewcommand{\footrulewidth}{0.4pt}
\fancyhead[L]{\small   \textit{Banaras Hindu University}}
\fancyhead[R]{\small   \textit{Botany Paper: BOB-602}}
\fancyfoot[C]{\small Page  \thepage\ of \pageref{LastPage}}


\newlist{parts}{enumerate}{2}
\setlist[parts,1]{label=(\alph*),leftmargin=*,itemsep=4pt,topsep=3pt}
\setlist[parts,2]{label=(\roman*),leftmargin=*,itemsep=2pt,topsep=2pt}

\newcommand{\pts}[1]{\hfill   \textit{[#1]}}


\begin{document}


\begin{center}
  {\large\bfseries BANARAS HINDU UNIVERSITY}\\[2pt]
  {\normalsize B.Sc. (Hons.) Semester VI Examination, 2013-14}\\[6pt]
  \rule{0.8\linewidth}{0.5pt}\\[6pt]
  {\large\bfseries Botany}\\[4pt]
  {\large\bfseries Paper: BOB-602 --- Microbiology and Plant Pathology}\\[4pt]
  {\normalsize Time: Three hours \hfill Full Marks: 70}
\end{center}

\rule{\linewidth}{0.4pt}

\smallskip



\noindent   \textit{Note: Attempt five questions in all, selecting two questions each from Section-A and B. Question No. 1 is compulsory. The figures in the right-hand margin indicate marks.}

\bigskip

\medskip


\noindent   \textbf{Question 1.} Answer the following in not more than 50 words each or fill in the blanks appropriately:\pts{$1  \times 10 = 10$}

\begin{parts}
  \item Define lag phase.
  \item Name two denitrifying bacteria.
  \item Name two lactic acid bacteria.
  \item Define T-DNA.
  \item Name a methanogenic bacterium.
  \item Name the pathogen of powdery mildew of Sheesham.
  \item Name two famines that were caused by plant pathogens.
  \item Define quarantine.
  \item Name a PR-protein.
  \item PUSA associated with IARI stands for \makebox[3cm]{\hrulefill}.
\end{parts}

\bigskip

\begin{center}
   \textbf{SECTION -- A}
\end{center}

\medskip


\noindent   \textbf{Question 2.} Give a detailed account of acetic acid production by different methods.\pts{15}

\medskip


\noindent   \textbf{Question 3.} Discuss critically the role of microbes in genetic engineering.\pts{15}

\medskip


\noindent   \textbf{Question 4.} Write explanatory notes on any two of the following:\pts{$7.5  \times 2 = 15$}

\begin{parts}
  \item A general account on Mycoplasma and Thermoplasma
  \item Retroviruses
  \item Mechanism of gene transfer in transduction
\end{parts}

\medskip


\noindent   \textbf{Question 5.} Write short notes on any three of the following:\dots\pts{$5  \times 3 = 15$}

\begin{parts}
  \item Bacterial growth curve
  \item Scope of Microbiology
  \item $ \text{N}_2$ fixing cyanobacteria
  \item Nitrification
\end{parts}

\bigskip

\begin{center}
   \textbf{SECTION -- B}
\end{center}

\medskip


\noindent   \textbf{Question 6.} How host physiology changes following attack by pathogens? Discuss the mechanisms involved with suitable examples.\pts{15}

\medskip


\noindent   \textbf{Question 7.} Give an account on Integrated Disease Management (IDM). Why is it considered as the best disease control strategy?\pts{15}

\medskip


\noindent   \textbf{Question 8.} Write explanatory notes on any two of the following:\pts{$7.5  \times 2 = 15$}

\begin{parts}
  \item Brief history of plant pathology and its scope
  \item Rust of linseed
  \item Wilt of tomato
\end{parts}

\medskip


\noindent   \textbf{Question 9.} Write notes on any three of the following:\pts{$5  \times 3 = 15$}

\begin{parts}
  \item Transmission and spread of plant diseases
  \item Smut of Bajra
  \item Mosaic of sugarcane
  \item Haustoria
\end{parts}

\vfill
\rule{\linewidth}{0.4pt}

\begin{center}\small
\href{https://bhu-exam.vercel.app/}{Click Here}\\[4pt]\href{https://me-aryan.vercel.app/}{Click Here}\\[4pt]\href{https://quashan-me.vercel.app/}{Click Here}
\end{center}

\end{document}
